\section{If AI fails}
\label{sec:bust}

The paper so far has asked what happens if AI succeeds and displaces labor. The opposite case
deserves the same treatment, and for most of this project it did not get it: displacement went
through a household engine, a fiscal module and thousands of balance sheets while the failure case
was assessed against historical analogies. That asymmetry flattered the failure case, so this
section runs it through the same machinery.

The answer is that an AI bust is a large fiscal event and a small credit event, which inverts the
displacement case exactly.

\subsection{What the AI side is financed with, in tiers}

Any claim about what a bust would do depends on how the investment was financed, so we set out
what is known in tiers, each with its own confidence label, and never blend them.
Table~\ref{tab:aitiers} gives them.

The hard tier is nine SEC registrants' own filings: \result{FilerCapexBn} billion dollars of
capital spending against \result{FilerOCFBn} billion of operating cash flow, a self-funding ratio
of \result{SelfFunding}, and \result{OnBalanceSheetAIDebtBn} billion dollars of on-balance-sheet
debt. On those books the debt-financed share of capital spending is
\result{DebtFinancedShare} percent. A second tier, verified from a Federal Reserve Bank source but
not from filings, is about \result{AIBankCommitmentsBn} billion dollars of large bank commitments
to AI-adjacent industries, of which roughly \result{AIDrawnBn} billion appears to be drawn. A third
tier, the equity value attributable to AI, has no defensible selection rule in our data and is
swept rather than asserted. The fourth tier, off-balance-sheet vehicles, GPU-backed lending, data
centre securitizations and vendor financing, is not sourced at all, and the Bank for International
Settlements describes exactly those structures as dominant in this market.

\begin{table}[htbp]
\centering
\caption{The AI side by tier of evidence. Tier 1 is measured from filings; tier 4 is the hole the
verdict turns on. Cells with no sourced value read n/a.}
\label{tab:aitiers}
\begin{tabular}{p{2.1cm}p{6.3cm}rp{3.1cm}}
\toprule
Tier & What it covers & Size (USD bn) & Confidence \\
\midrule
1 HARD & on-balance-sheet AI debt, nine named SEC filers & 458.4 & hard \\
2 VERIFIED SECONDARY & large bank C and I commitments to AI-adjacent industries & 450.0 & verified secondary \\
2b VERIFIED SECONDARY, drawn & the drawn share of those commitments, implied by the same source & 162.0 & implied from outstanding 9 pct vs committed 25 pct of tier 1 \\
3 SCENARIO & listed AI-exposed equity, swept share of US nonfinancial corporate equity & 7,200 to 21,598 & SCENARIO \\
4 NOT SOURCED & off-balance-sheet data centre vehicles, GPU-backed loans, private credit to data centres, data centre ABS and CMBS, vendor financing & n/a & NOT SOURCED \\
MEMO & chip supply chain & n/a & named \\
\bottomrule
\end{tabular}

\end{table}

\subsection{The verdict on the financing structure, which is conditional}

There are two historical patterns for an asset bust of this kind, and they have very different
fiscal consequences. In the equity-led case of 2000 to 2002, federal receipts fell
\result{EquityLedReceiptsPct} percent and corporate tax \result{EquityLedCorpPct} percent. In the
debt-led case of 2007 to 2009, receipts fell \result{CreditLedReceiptsPct} percent and corporate
tax \result{CreditLedCorpPct} percent. Both are verified from the published aggregates.

On the nine filers' own books, at \result{DebtFinancedShare} percent debt-financed, the structure
sits near the equity-led pattern. That is where a paper could stop, and it would be wrong to. The
verdict flips out of that pattern on \result{FlipToBetweenBn} billion dollars of additional
debt-financed capital spending, which is \result{FlipToBetweenShare} percent of the bank
commitments already identified in tier 2, and reaches the debt-led pattern on
\result{FlipToDebtLedBn} billion, \result{FlipToDebtLedShare} percent of them. The drawn portion of
those commitments alone is already larger than the first threshold. Tier 4 is unmeasured.

\begin{quote}
The defensible statement is that on the nine filers' own balance sheets the structure resembles the
equity-led case. The statement that the structure resembles the equity-led case is not defensible,
and we do not make it. The verdict is conditional on a quantity we have not measured, and the
direction of the unmeasured quantity is towards the worse case.
\end{quote}

\subsection{The bust run through the same engine}

Converted to a wage-bill equivalent, an AI bust costs \result{BustWageLow} to
\result{BustWageHigh} percent of the wage bill, which is \result{BustRatioLow} to
\result{BustRatioHigh} times smaller than the ten percent displacement case this paper has used
throughout. Household credit losses are \result{BustCreditLow} to \result{BustCreditHigh} billion
dollars, against tens of billions in the displacement case. Output falls \result{BustGDPLow} to
\result{BustGDPHigh} percent, well inside the fall in the Federal Reserve's severely adverse
scenario. All of these are scenario figures.

Federal receipts, on the other hand, fall \result{BustReceiptsLow} to \result{BustReceiptsHigh}
billion dollars on the verified historical episodes. That is an order of magnitude larger than the
terminal-year fiscal loss at ten percent displacement, and it arrives through a completely
different door: capital gains and corporate receipts rather than wage taxes.

Why the bust transmits so weakly to households is measured rather than assumed, and it is the same
ownership fact that drove Section~\ref{sec:fiscal}. Corporate equity is extraordinarily
concentrated: the top one percent of households hold \result{TopOnePctEquity} percent of it and the
bottom half hold \result{BottomHalfEquity} percent. A fall in equity value therefore lands on the
households with the lowest propensity to consume out of wealth, measured at \result{MPCWealth}
cents in the dollar per year. The same concentration that makes the capital tax base hard to reach
makes the wealth channel weak.

Two caveats belong with these numbers. The propensity to consume out of wealth is a sourced
parameter checked against two episodes, not a parameter we estimated. And the path through which
output falls does not feed bank-channel damage back into itself, so the output fall is understated
by an amount we have not sized.

\subsection{The payoff by holder, and what a holdings table cannot show}

Table~\ref{tab:payoff} scores each holder by its share of the wage side and its share of the AI
side, and reports the net position in three outcomes. Only the federal government and banks lose in
every column. Everyone else holds enough of the AI side to offset what they lose on the wage side
if AI succeeds: households in particular hold \result{HouseholdsAISide} percent of the AI side
against six percent of the wage side, and gain \result{HouseholdPayoffSuccess} under success where
the federal government loses \result{FedPayoffSuccess}.

\begin{table}[htbp]
\centering
\caption{The position of each holder across three outcomes, scored on holdings. The wage side is
measured; the AI side is scenario and a lower bound.}
\label{tab:payoff}
\begin{tabular}{lrrrrr}
\toprule
Holder & Wage side & AI side & AI fails & Partial & AI succeeds \\
\midrule
Federal government & 0.3215 & 0.0100 & -0.0132 & -0.1657 & -0.3115 \\
Banks & 0.1844 & 0.0318 & -0.0337 & -0.1081 & -0.1525 \\
Rest of the world & 0.1801 & 0.1808 & -0.1826 & -0.1804 & +0.0007 \\
Other financial & 0.1532 & 0.1633 & -0.1648 & -0.1583 & +0.0101 \\
Households & 0.0607 & 0.3842 & -0.3848 & -0.2224 & +0.3235 \\
State and local & 0.0462 & 0.0363 & -0.0367 & -0.0413 & -0.0100 \\
Insurers & 0.0203 & 0.0207 & -0.0209 & -0.0205 & +0.0003 \\
Unallocated remainder & 0.0164 & 0.0979 & -0.0981 & -0.0572 & +0.0816 \\
Pension funds & 0.0110 & 0.0437 & -0.0439 & -0.0274 & +0.0327 \\
Nonfinancial business & 0.0062 & 0.0313 & -0.0313 & -0.0187 & +0.0250 \\
\bottomrule
\end{tabular}

\end{table}

The table understates the federal position in the first column and the reason matters more than the
table. It scores holders by what they hold, so a government holding one percent of the AI side
scores as barely exposed when AI fails, at \result{FedPayoffFails}. But federal receipts fall
\result{BustReceiptsLow} to \result{BustReceiptsHigh} billion dollars in exactly that case. The
state's claim on the AI upside is fiscal rather than proprietary: it is the capital tax of
Section~\ref{sec:fiscal}, not a portfolio. A holdings table cannot see a tax claim, which is why
this paragraph sits beside the table rather than in a footnote.

Read that way, the two facts are one fact. Being barely exposed to an AI bust, in holdings terms,
and being unhedged on the AI upside are the same statement about a government whose only claim on
AI capital is a tax that reaches \result{TauKBarkai} percent of the surplus.

\subsection{Two directions, two tax bases}

The instruments that would protect against the two failure directions are nearly disjoint, but not
for the reason it would be natural to assume. It is not that a bust leaves the wage side untouched:
it touches it, and the size of that is reported above. It is that the bust reaches the wage side
through a channel one to two orders of magnitude smaller than displacement does, while reaching the
federal budget through capital gains and corporate receipts, which displacement barely touches.

The two failure directions therefore hit the same institution through different tax bases.
Instruments that protect the wage tax base do not protect the capital tax base, and the reverse.
That is the organising conclusion this paper carries into
Section~\ref{sec:institutions}: one capital tax rate is simultaneously the state's only claim on
the AI upside and the thing that would have to rise to close the fiscal condition after
displacement. It is being asked to do two opposite jobs.
